% FILE: 06_contribution_to_thesis.tex
% Project: prompts — Downstream Manufacturing Prompt Corpus

\section{Contribution to the Dissertation}
\label{sec:prompts-contribution}

\subsection{Contribution to Prediction vs.~Coordination}
\label{sec:prompts-contribution-prediction}

The dissertation's central claim is that the failure mode of AI-adjacent systems is
not inaccurate prediction but incoherent coordination---the inability of components to
prove to one another that a shared invariant has been maintained. The prompt corpus
contributes to this claim by demonstrating what coordination without proof looks like
in practice and then specifying what proof-bearing coordination requires.

The old regime of downstream engineering coordination is represented by informal
specifications: a directive that says ``implement Inductive Miner'' without specifying
the return type, the authority source, the completion criterion, or the blocking
claims. Such a directive cannot coordinate because a downstream engineer has no
mechanism to verify that the implementation satisfies the research program's
requirements. The engineer might implement an algorithm that produces structurally
correct-looking output that is unsound, and no artifact would capture the failure
before it propagates into M\&A claims.

The corpus's contribution is the proof-bearing directive format: each directive
carries a named law type for every refusal, a compile-time fixture for every
structural guarantee, an authority source for every mandate, and an explicit list of
blocking claims that must carry \texttt{[BLOCKED: GAP\_NNN]} labels until the
underlying gap is closed. This format makes coordination verifiable: the directive
either satisfies its completion criterion or it does not, and the determination is
mechanical.

\subsection{Contribution to Process Evidence}
\label{sec:prompts-contribution-evidence}

The corpus makes a concrete contribution to the process evidence chain by specifying
the \texttt{ProcessIntelligenceVerificationReceipt} schema as the atomic unit of
board-admissible evidence. This schema links three artifacts that are conventionally
kept separate: the source event log (captured by \texttt{target\_log\_hash}), the
process model (captured by \texttt{process\_model\_hash}), and the conformance result
(captured by \texttt{fitness}, \texttt{precision}, and \texttt{throughput\_days}),
all bound together by an Ed25519 signature from a registered validator.

The fitness equation $f(\sigma, N) = \frac{1}{2}(1 - m/c) + \frac{1}{2}(1 - r/p)$
(where $m$ is missing tokens, $c$ is consumed tokens, $r$ is remaining tokens, and
$p$ is produced tokens) and the A* alignment cost function $K(a, \gg) = 1$,
$K(\gg, t) = 1$, $K(a, t) = 0$ iff $\operatorname{label}(t) = a$ are embedded as
requirements in the refactor mandate. These formulas establish that conformance
checking must be alignment-based, not token-replay-based, because token replay
produces fitness values that can be artificially inflated by model simplification.
The corpus thus contributes to the process evidence literature by mandating the
stronger conformance method.

\subsection{Contribution to Manufacturing Architecture}
\label{sec:prompts-contribution-manufacturing}

The corpus contributes a formal manufacturing architecture for the translation from
research finding to engineering specification. The architecture has four layers:

\begin{enumerate}
  \item \textbf{Research authority layer}: Gap register, adversarial red-team
    findings, conformance authority maps, M\&A claim taxonomies. These authorize
    directives but do not themselves specify engineering work.
  \item \textbf{Directive layer} (the corpus): Translates research authority into
    sequenced, gated engineering tasks with explicit completion criteria and blocking
    claims.
  \item \textbf{Blueprint layer} (execution plans): Provides mathematical and
    cryptographic detail for the most complex directive requirements---lattice
    mathematics, BLAKE3/Ed25519 receipt construction, typestate enforcement patterns,
    absence-proof fixture strategies.
  \item \textbf{Verification layer}: Compile-pass and compile-fail fixtures, receipt
    schemas, benchmark targets, and gap-register updates that confirm completion.
\end{enumerate}

This four-layer architecture is the manufacturing pipeline for process intelligence
engineering. It demonstrates that a research program can maintain formal authorization
chains through all layers of translation from theorem to implementation.

\subsection{Contribution to Capital-Flow Infrastructure}
\label{sec:prompts-contribution-capital}

The corpus makes its most novel contribution at the intersection of process evidence
and capital-flow claims. The ggen integration directive specifies a complete pipeline
from receipt-shaped process evidence to board-admissible M\&A claims, with EBITDA
impact percentages tied to alignment fitness thresholds:

\begin{itemize}
  \item Fitness $\geq 0.95$: Conformance claim with \texttt{HIGH} confidence,
    $+5\%$ EBITDA impact, automatic board admissibility.
  \item $0.85 \leq \text{fitness} < 0.95$: Qualified conformance with \texttt{MEDIUM}
    confidence, $+3\%$ EBITDA impact, board override required.
  \item Fitness $< 0.85$: Non-conformance detection, \texttt{LOW} confidence, claim
    refused, risk mitigation emitted with $-8\%$ to $-15\%$ EBITDA impact.
\end{itemize}

The EBITDA figures are derived from the M\&A board claim taxonomy documents
(\texttt{ma/define\_board\_claim\_taxonomy.md}), which are themselves products of the
research program. Whether these figures are grounded in empirical studies or in the
research program's own assertion is an open question (see
Section~\ref{sec:prompts-open-questions}). The contribution, regardless, is formal:
it establishes that capital-flow claims must carry receipt chains, and that the
receipt chain is the legally admissible foundation for acquisition valuation
adjustments.
